% !TEX root = main.tex
% =====================================================================
%  Presentation du projet POO : moteur masses-ressorts (Phelma 3A SEI)
%  Charte reprise de la soutenance de stage (theme << ressorts >>, dans theme/).
%
%  Compilation :  make            -> build/main.pdf
%                 make fiche      -> build/main-fiche.pdf (fiche d'orateur)
%                 make captures   -> regenere figures/scenes/*.png depuis le
%                                    programme (cible CMake << capture >>)
% =====================================================================
\documentclass[aspectratio=169,11pt]{beamer}

\ifdefined\astfichemode
  \usepackage{pgfpages}
  \setbeameroption{show only notes}
  \pgfpagesuselayout{4 on 1}[a4paper,landscape,border shrink=4mm]
\fi

\makeatletter
\def\input@path{{theme/}}
\makeatother
\usetheme[progressbar=true,sectionpages=true]{ressorts}

% !TEX root = main.tex
% Preambule : volontairement minimal (beamer charge deja hyperref).
\usepackage[utf8]{inputenc}
\usepackage[french]{babel}

% La bibliotheque babel de TikZ neutralise la ponctuation active (; : ! ?)
% que babel-french installe : sans elle, tout chemin TikZ casse.
\usetikzlibrary{calc,positioning,arrows.meta,fit,backgrounds,patterns,
  decorations.pathreplacing,decorations.pathmorphing,shapes.geometric,babel}

\usepackage{booktabs}
\usepackage{array}
\usepackage{fontawesome5}
\usepackage{appendixnumberbeamer}

\usepackage{listings}
\lstdefinestyle{cpp}{
  language=C++,
  basicstyle=\ttfamily\scriptsize,
  keywordstyle=\color{astblue}\bfseries,
  commentstyle=\color{textgrey}\itshape,
  stringstyle=\color{datagreen!60!black},
  emph={Particle,Entity,World,Spring,Force,Drawable,Vec2,Integrator,Euler,Verlet,
        Gravity,Drag,Wind,Renderer,ParticleList,Scene,SceneInstance},
  emphstyle=\color{phelmared},
  morekeywords={override,nullptr,constexpr,auto},
  breaklines=true,
  frame=single,
  framerule=0.4pt,
  rulecolor=\color{astrulegrey},
  backgroundcolor=\color{astframefill},
  xleftmargin=2mm,xrightmargin=2mm,
  showstringspaces=false,
  tabsize=2,
  literate={é}{{\'e}}1 {è}{{\`e}}1 {à}{{\`a}}1,
}
\lstset{style=cpp}

\pdfsuppresswarningpagegroup=1
\hypersetup{
  pdftitle={Moteur masses-ressorts : projet de programmation orientée objet},
  pdfauthor={Léo Dorval, Nicolas Marcel, Audric Siguret},
  pdfsubject={Projet POO C++ - Phelma 3A SEI},
  pdfpagemode=UseNone,
  pdfstartview=Fit,
  pdfdisplaydoctitle=true,
  hidelinks,
}

% !TEX root = main.tex
% Outils de mise en page : reprise de ceux de la soutenance de stage
% (diagramoverlay, callouts, pastilles, fiche d'orateur), plus quelques
% outils propres au projet (captures, classes UML, puces de fichier).

% Hauteur utile pour une image : la page moins les deux bandeaux, moins un jeu.
\newlength{\astdiagheight}
\setlength{\astdiagheight}{\dimexpr\paperheight-\astheadheight-\astfootheight-0.7cm\relax}

% Separateur entre mots-cles :  surface \astsep fuite
% \texorpdfstring : dans les metadonnees du PDF (\author -> pdfauthor), une
% simple virgule, sinon hyperref avertit << Token not allowed in a PDF string >>
\newcommand{\astsep}{\texorpdfstring{\;\textperiodcentered\;}{, }}

% ---------------------------------------------------------------------------
% diagramoverlay : on dessine sur une image en coordonnees relatives,
% (0,0) = coin bas gauche, (1,1) = coin haut droit.
% ---------------------------------------------------------------------------
\newif\ifastdiaggrid
\newcommand{\astdrawgrid}{%
  \draw[cyan!60,very thin,step=0.1] (0,0) grid (1,1);
  \foreach \i in {0,0.1,...,1.01}{%
    \node[font=\tiny,text=cyan!70!black,anchor=north] at (\i,0)
      {\pgfmathprintnumber[fixed,precision=1]{\i}};
    \node[font=\tiny,text=cyan!70!black,anchor=east] at (0,\i)
      {\pgfmathprintnumber[fixed,precision=1]{\i}};}}

\NewDocumentEnvironment{diagramoverlay}{O{} m}{%
  \begin{center}%
  \begin{tikzpicture}
    \node[anchor=south west,inner sep=0,outer sep=0] (img) at (0,0)
      {\includegraphics[width=\textwidth,height=\astdiagheight,keepaspectratio,#1]{#2}};
    \begin{scope}[x={(img.south east)},y={(img.north west)}]
      \ifastdiaggrid\astdrawgrid\fi
}{%
    \end{scope}
  \end{tikzpicture}%
  \end{center}%
}

\tikzset{astcallout/.style={draw=phelmared,fill=white,text=inpnavy,rounded corners=1.5pt,
  line width=0.8pt,font=\footnotesize,inner sep=3pt,align=left}}
\NewDocumentCommand{\highlight}{D<>{1-} O{phelmared} m m m m}{%
  \only<#1>{\draw[#2,line width=1.2pt,rounded corners=2pt] (#3,#4) rectangle (#5,#6);}}
\NewDocumentCommand{\spotlight}{D<>{1-} m m m m}{%
  \only<#1>{\fill[white,fill opacity=0.66,even odd rule]
    (0,0) rectangle (1,1) (#2,#3) rectangle (#4,#5);}}
\tikzset{astnum/.style={circle,fill=phelmared,draw=white,line width=0.6pt,text=white,
  font=\scriptsize\bfseries,inner sep=0pt,minimum size=4.2mm}}
\newcommand{\astmark}[3]{\node[astnum] at (#1,#2) {#3};}
\NewDocumentCommand{\callout}{D<>{1-} O{} m m m}{%
  \only<#1>{\node[astcallout,#2] at (#3,#4) {#5};}}
% Etiquette reliee par une fleche :  \astpoint{x}{y}{xe}{ye}{ancre}{texte}
\newcommand{\astpoint}[6]{%
  \node[astcallout,anchor=#5] (astpt) at (#3,#4) {#6};
  \draw[-{Latex[length=1.6mm]},phelmared,line width=0.8pt] (astpt.#5) -- (#1,#2);}
\newcommand{\astpill}[1]{\tikz[baseline=(astp.base)]\node[astnum] (astp) {#1};}

% ---------------------------------------------------------------------------
% Projet : puces de code et de fichier
% ---------------------------------------------------------------------------
% \code{dynamic_cast}   mot de code dans le texte
\newcommand{\code}[1]{{\ttfamily\color{combislate}#1}}
% \filechip{world.hpp}  petite etiquette de fichier, grise
\newcommand{\filechip}[1]{\tikz[baseline=(fc.base)]\node[fill=asttrack,text=textgrey,
  rounded corners=1pt,inner xsep=2pt,inner ysep=1.2pt,font=\ttfamily\tiny] (fc) {#1};}

% ---------------------------------------------------------------------------
% Classes UML dans le langage des schemas : bord colore, fond pale.
%   \node[cls=memborder/memfill] (P) {\clsname{Particle}\clsbody{pos, speed\\mass}};
% ---------------------------------------------------------------------------
\tikzset{
  cls/.style args={#1/#2}{draw=#1,fill=#2,line width=0.8pt,rounded corners=1.6mm,
    align=center,inner sep=1.4mm,font=\scriptsize,text=astbodyfg},
  abstract/.style={dash pattern=on 3pt off 1.5pt},
  blackbox/.style={draw=floorgrey,fill=asttrack,line width=0.8pt,rounded corners=1.6mm,
    align=center,inner sep=1.4mm,font=\scriptsize,text=textgrey},
  inherit/.style={-{Triangle[open,length=2.4mm,width=2.4mm]},line width=0.7pt,draw=combislate},
  vinherit/.style={inherit,draw=phelmared,line width=1pt},
  owns/.style={{Diamond[length=2.6mm,width=1.8mm]}-{Latex[length=1.6mm]},line width=0.7pt,draw=combislate},
  uses/.style={-{Latex[length=1.6mm]},line width=0.7pt,draw=combislate,dash pattern=on 2pt off 1.5pt},
  flow/.style={-{Latex[length=2mm]},line width=1pt,draw=inpnavy},
}
\newcommand{\clsname}[1]{{\bfseries\small #1}}
\newcommand{\clsabstract}[1]{{\bfseries\itshape\small #1}}
% Le corps passe par un tabular : dans un noeud TikZ, un \\ remet la police
% du noeud, et un \\ dans un groupe casse le noeud.
\newcommand{\clsbody}[1]{\\[-0.2ex]{\color{textgrey}\tiny
  \renewcommand{\arraystretch}{0.9}\begin{tabular}{@{}c@{}}#1\end{tabular}}}
\newcommand{\clsleaf}[1]{{\bfseries\footnotesize #1}}

% ---------------------------------------------------------------------------
% \scenecard{fichier}{nom}{hauteur}  capture d'une scene dans un cadre
% d'etage pointille, avec le nom sur le bord bas (comme les etages des schemas)
% ---------------------------------------------------------------------------
\newcommand{\scenecard}[3]{%
  \begin{tikzpicture}
    \node[inner sep=0] (sc) {\includegraphics[height=#3]{#1}};
    \draw[floorgrey,line width=0.8pt,dash pattern=on 3pt off 2pt,rounded corners=1.5mm]
      ($(sc.south west)+(-1mm,-1mm)$) rectangle ($(sc.north east)+(1mm,1mm)$);
    \node[fill=white,draw=combislate,line width=0.6pt,font=\tiny\bfseries,
      inner xsep=1.5mm,inner ysep=0.8mm] at ($(sc.south)+(0,-1mm)$) {\MakeUppercase{#2}};
  \end{tikzpicture}}

% ---------------------------------------------------------------------------
% \astbig[couleur]{400}{particules}   grand nombre et sa legende
% ---------------------------------------------------------------------------
\NewDocumentCommand{\astbig}{O{inpnavy} m m}{%
  \begin{center}
    {\fontsize{24}{28}\selectfont\bfseries\color{#1}#2}\\[0.3mm]
    {\footnotesize\color{textgrey}#3}
  \end{center}}

% ---------------------------------------------------------------------------
% Fiche d'orateur (cible << fiche >>) : pour chaque diapositive, quelques
% reperes courts, a imprimer et a garder a cote.
%   \astcue{ \cue{idee, \cuekey{mot cle}} \cuewarn{piege} \cuenext{transition} }
% ---------------------------------------------------------------------------
\newcommand{\cuekey}[1]{\textbf{\color{phelmared}#1}}
\newcommand{\cue}[1]{\item[{\color{phelmared}$\rightarrow$}] #1}
\newcommand{\cuewarn}[1]{\item[{\color{dummyorange}\footnotesize\faExclamationTriangle}] {\color{dummyorange!80!black}#1}}
\newcommand{\cuenext}[1]{\item[{\color{astblue}\footnotesize\faShare}] {\color{astblue}\itshape #1}}
\ifdefined\astfichemode
  \setbeamertemplate{note page}{%
    \vspace*{2mm}%
    \begin{minipage}[t]{0.36\linewidth}\vspace{0pt}%
      \fbox{\insertslideintonotes{0.30}}%
    \end{minipage}\hspace{3mm}%
    \begin{minipage}[t]{0.58\linewidth}\vspace{1mm}%
      {\normalsize\bfseries\color{inpnavy}\insertframetitle}\par\vspace{1mm}%
      {\scriptsize\color{textgrey}diapositive \insertframenumber}%
    \end{minipage}\par\vspace{2mm}%
    \begin{minipage}[t]{0.97\linewidth}%
      \small\raggedright\insertnote
    \end{minipage}}
  \newcommand{\astcue}[1]{\note{\begin{itemize}\setlength{\itemsep}{1.2mm}#1\end{itemize}}}
\else
  \newcommand{\astcue}[1]{}
\fi

% ---------------------------------------------------------------------------
% Fil rouge de la partie 2 : les huit etapes d'un pas de simulation, en bas
% de la diapositive. L'etape courante est pleine, les suivantes estompees.
%   \fil{4}   a placer dans la frame (dessin en surimpression)
% ---------------------------------------------------------------------------
\newcommand{\filsteps}{Scène,Particules,World,Forces,Intégration,Contraintes,Dessin,Souris}
\newcommand{\fil}[1]{%
  \begin{tikzpicture}[remember picture,overlay]
    \begin{scope}[shift={($(current page.south)+(0,\astfootheight+1.6mm)$)}]
      \foreach \lab [count=\k] in \filsteps {
        \pgfmathsetmacro{\fx}{(\k-4.5)*1.72}
        \ifnum\k<#1\relax
          \node[draw=inpnavy,fill=white,text=inpnavy,rounded corners=1pt,font=\tiny,
            minimum width=1.55cm,inner ysep=1pt] at (\fx,0) {\k\enspace\lab};
        \else\ifnum\k=#1\relax
          \node[draw=phelmared,fill=phelmared,text=white,rounded corners=1pt,font=\tiny\bfseries,
            minimum width=1.55cm,inner ysep=1pt] at (\fx,0) {\k\enspace\lab};
        \else
          \node[draw=astmutedbadge,fill=white,text=floorgrey,rounded corners=1pt,font=\tiny,
            minimum width=1.55cm,inner ysep=1pt] at (\fx,0) {\k\enspace\lab};
        \fi\fi
      }
    \end{scope}
  \end{tikzpicture}}

% ---------------------------------------------------------------------------
% Schemas draw.io : sources dans schemas/*.drawio, PDF exportes par
% << make schemas >> dans figures/schemas/. Ils sont dessines a l'echelle 2 :
% inclus a 0,5, leurs textes ont la taille des textes des diapositives
% (21 px = \scriptsize). Ils ne sont reduits que s'ils depassent la place.
%   \schema{hierarchie}          hauteur max : toute la zone utile
%   \schema[5.3cm]{forces}       hauteur max donnee (diapos avec fil rouge)
% ---------------------------------------------------------------------------
\NewDocumentCommand{\schema}{O{\astdiagheight} m}{%
  \adjustbox{scale=0.5,max width=\linewidth,max totalheight=#1}%
    {\includegraphics{figures/schemas/#2.pdf}}}


% \astdiaggridtrue   % <- decommenter pour calibrer les overlays

\title[Moteur masses-ressorts]{Moteur masses-ressorts en C++}
\subtitle{Projet de programmation orientée objet}
\author[Dorval\astsep Marcel\astsep Siguret]{Léo DORVAL\astsep Nicolas MARCEL\astsep Audric SIGURET}
\institute[Phelma]{Grenoble INP -- Phelma\astsep 3\textsuperscript{e} année SEI}
\date{2026}
\astinternship{Programmation orientée objet en C++\astsep novembre 2026}

\begin{document}
% !TEX root = ../main.tex
% Ouverture : titre et ce que fait le moteur

% ======================================================================= titre
\begin{frame}[plain]
  \titlepage
\end{frame}

% ============================================================= ce que fait le moteur
\begin{frame}{Ce que fait le moteur}
  \framesubtitle{des particules, des ressorts, des forces}
  \begin{columns}[c,onlytextwidth]
    \begin{column}{0.58\textwidth}
      \begin{diagramoverlay}[width=\linewidth]{figures/scenes/scene_05.png}
        \astpoint{0.737}{0.94}{0.79}{0.80}{west}{ancre\\épinglée}
        \astpoint{0.722}{0.78}{0.79}{0.60}{west}{ressort\\étiré}
        \astpoint{0.5}{0.32}{0.5}{0.10}{north}{400 particules\astsep 1\,482 ressorts}
      \end{diagramoverlay}
    \end{column}
    \begin{column}{0.38\textwidth}
      \astbig{13}{scènes, du pendule au tissu}
      \vspace{1mm}
      \astbig[phelmared]{4}{forces : poids, frottement, ressort, vent}
      \vspace{1mm}
      \astbig[timagreen]{1/120\,s}{pas fixe, Euler ou Verlet en direct}
    \end{column}
  \end{columns}
  \astcue{
    \cue{montrer le tissu en direct, \cuekey{attraper une particule}}
    \cuenext{comment c'est organisé}
  }
\end{frame}
          % titre, ce que fait le moteur
% !TEX root = ../main.tex
% Partie 1 : architecture (dossiers, hierarchie des classes)

\section{Architecture}

% ================================================================ architecture
\begin{frame}{Quatre dossiers, une frontière}
  \framesubtitle{le moteur ne connaît que l'interface \code{Renderer}}
  \centering
  \begin{tikzpicture}[x=1cm,y=1cm]
    % physics : la bibliotheque
    \node[draw=floorgrey,line width=0.8pt,dash pattern=on 3pt off 2pt,rounded corners=3mm,
      fill=floorfill,minimum width=9.2cm,minimum height=3.3cm] (phys) at (4.6,1.65) {};
    \node[fill=white,draw=combislate,line width=0.6pt,font=\scriptsize\bfseries,
      inner xsep=2mm,inner ysep=1mm] at (phys.south) {PHYSICS/ \astsep{} LIBPHYSICS.A};
    \node[cls=datagreen/datafill,minimum width=1.9cm] (w) at (1.5,2.35) {\clsname{World}\clsbody{possède tout}};
    \node[cls=datagreen/datafill,minimum width=1.9cm] (p) at (3.7,2.35) {\clsname{Particle}\clsbody{position, masse}};
    \node[cls=datagreen/datafill,minimum width=1.9cm] (f) at (5.9,2.35) {\clsname{forces/}\clsbody{poids, ressort, vent}};
    \node[cls=datagreen/datafill,minimum width=1.9cm] (i) at (8.1,2.35) {\clsname{integrators/}\clsbody{Euler, Verlet}};
    \node[cls=memborder/memfill,abstract,minimum width=4cm] (r) at (4.6,0.95)
      {\clsabstract{Renderer}\clsbody{drawCircle, drawLine : virtuelles pures}};
    % scenes et app au-dessus
    \node[cls=dummyorange/dummyfill,minimum width=3.2cm] (sc) at (2.4,4.35)
      {\clsname{scenes/}\clsbody{13 scènes = 13 fonctions}};
    \node[cls=ctrlred/ctrlfill,minimum width=3.2cm] (app) at (6.8,4.35)
      {\clsname{app/main.cpp}\clsbody{boucle, clavier, souris}};
    % render : boite noire
    \node[blackbox,minimum width=3.0cm,minimum height=1.3cm] (sf) at (12.0,0.95)
      {{\bfseries\small render/}\\[0.3ex]SFMLRender\\{\tiny boîte noire}};
    \node[blackbox,minimum width=3.0cm] (lib) at (12.0,2.75) {SFML 2.6\\{\tiny FetchContent, statique}};
    % liens
    \draw[uses] (sc.south) -- (sc.south |- phys.north);
    \draw[uses] (app.south) -- (app.south |- phys.north);
    \draw[uses] (app.west) -- (sc.east);
    \draw[inherit] (sf.west) -- node[above,font=\tiny,text=textgrey]{implémente} (r.east);
    \draw[uses] (sf.north) -- (lib.south);
    \draw[uses] (app.east) -| (sf.north east);
    % la frontiere
    \draw[phelmared,line width=1.2pt,dash pattern=on 5pt off 3pt] (9.95,-0.25) -- (9.95,5.0);
    \node[font=\tiny\bfseries,text=phelmared,anchor=north,align=center] at (9.95,-0.3)
      {aucun en-tête SFML à gauche};
  \end{tikzpicture}
  \astcue{
    \cue{\cuekey{physics} compile en \code{.a}, sans SFML : vérifié avec \code{grep}}
    \cue{le rendu SFML : on n'en parle pas, c'est une boîte noire}
    \cuenext{ce qu'il y a dans physics}
  }
\end{frame}

% ================================================================ hierarchie
\begin{frame}{La hiérarchie des classes}
  \framesubtitle{une seule \code{Entity} partagée, grâce à l'héritage virtuel}
  \centering
  \begin{tikzpicture}[x=1cm,y=1cm]
    % racine
    \node[cls=combislate/white,minimum width=2.6cm] (E) at (5.6,5.75)
      {\clsname{Entity}\clsbody{id, name\astsep \code{static} current\_ids}};
    % interfaces
    \node[cls=memborder/memfill,abstract,minimum width=2.4cm] (D) at (3.0,3.95)
      {\clsabstract{Drawable}\clsbody{couleurs\\\code{draw(Renderer\&) = 0}}};
    \node[cls=memborder/memfill,abstract,minimum width=2.4cm] (F) at (8.2,3.95)
      {\clsabstract{Force}\clsbody{\code{compute(const Particle\&) = 0}}};
    % concret
    \node[cls=datagreen/datafill,minimum width=2.2cm] (P) at (0.95,2.05)
      {\clsname{Particle}\clsbody{pos, speed, mass\\forces : \code{Force\&}}};
    \node[cls=datagreen/datafill,minimum width=2.2cm] (S) at (5.6,2.05)
      {\clsname{Spring}\clsbody{$L_0$, $k$, $\rho$\\2 \code{Particle*}}};
    \node[cls=datagreen/datafill] (G) at (7.6,2.05) {\clsleaf{Gravity}\clsbody{$m\vec g$}};
    \node[cls=datagreen/datafill] (Dr) at (8.75,2.05) {\clsleaf{Drag}\clsbody{$-b\vec v$}};
    \node[cls=datagreen/datafill] (Wi) at (9.95,2.05) {\clsleaf{Wind}\clsbody{lambda $w(t)$}};
    % integrateurs
    \node[cls=memborder/memfill,abstract,minimum width=2.4cm] (I) at (12.9,3.95)
      {\clsabstract{Integrator}\clsbody{\code{step(particles, dt) = 0}}};
    \node[cls=datagreen/datafill] (Eu) at (12.15,2.05) {\clsleaf{Euler}\clsbody{semi-implicite}};
    \node[cls=datagreen/datafill] (Ve) at (13.65,2.05) {\clsleaf{Verlet}\clsbody{garde $x_{n-1}$}};
    % monde
    \node[cls=ctrlred/ctrlfill,minimum width=3.4cm] (W) at (5.6,0.2)
      {\clsname{World}\clsbody{copie interdite : \code{= delete}}};
    % heritages virtuels (le diamant)
    \draw[vinherit] (P.north) |- (E.west);
    \draw[vinherit,-] (D.north) -- (D.north |- E.west);
    \draw[vinherit] (F.north) |- (E.east);
    \draw[inherit] ([xshift=-5mm]S.north) |- (D.east);
    \draw[inherit] ([xshift=5mm]S.north) |- (F.west);
    \node[font=\tiny\bfseries,text=phelmared] at (5.6,3.95) {diamant};
    % heritages simples
    \draw[inherit] (P.east) -| ([xshift=-3mm]D.south);
    \draw[inherit] (G.north) -- (G.north |- F.south);
    \draw[inherit] (Dr.north) -- (Dr.north |- F.south);
    \draw[inherit] (Wi.north) |- (F.east);
    \draw[inherit] (Eu.north) -- (Eu.north |- I.south);
    \draw[inherit] (Ve.north) -- (Ve.north |- I.south);
    % possession
    \draw[owns] (W.west) -| node[pos=0.8,left,font=\tiny,text=textgrey,align=right]{\code{deque}\\\code{<Particle>}} (P.south);
    \draw[owns] (W.north) -- node[right,font=\tiny,text=textgrey]{\code{unique\_ptr<Spring>}} (S.south);
    \draw[owns] (W.east) -- node[above,pos=0.45,font=\tiny,text=textgrey]{\code{unique\_ptr<Integrator>}} (10.95,0.2) |- (I.west);
    % legende
    \begin{scope}[shift={(11.6,0.5)},font=\tiny,text=textgrey]
      \draw[vinherit] (0,0.45) -- (0.6,0.45); \node[anchor=west] at (0.7,0.45) {héritage virtuel};
      \draw[inherit] (0,0.1) -- (0.6,0.1); \node[anchor=west] at (0.7,0.1) {héritage};
      \draw[owns] (0,-0.25) -- (0.6,-0.25); \node[anchor=west] at (0.7,-0.25) {possède};
    \end{scope}
  \end{tikzpicture}
  \astcue{
    \cue{\cuekey{Spring} est à la fois une force et un dessin}
    \cue{sans \code{virtual} : deux \code{Entity}, deux identifiants}
    \cuenext{comment ces classes travaillent ensemble à chaque pas}
  }
\end{frame}
   % dossiers et frontiere, hierarchie des classes
% !TEX root = ../main.tex
% Partie 2 : on suit un pas de simulation de bout en bout. Chaque notion du
% cours apparait au moment ou on la rencontre. Le fil rouge (\fil{n}) en bas
% de chaque diapositive rappelle ou l'on est.

\section{Un pas de simulation}

% ============================================================ la carte
\begin{frame}{Le parcours d'un pas}
  \framesubtitle{huit étapes, de la scène à la souris}
  \centering
  \schema{parcours}
  \astcue{
    \cue{la \cuekey{carte} de toute la partie : on la suit étape par étape}
    \cue{1 à 3 : une seule fois \astsep{} 4 à 6 : deux fois par image \astsep{} 7 : une fois par image}
    \cue{pas fixe : \cuekey{Verlet suppose un dt constant}}
    \cuenext{étape 1 : d'où vient une scène}
  }
\end{frame}

% ============================================================ 1. scene
\begin{frame}{Une scène est une fonction}
  \framesubtitle{\code{std::function}, lambdas et un \code{template}}
  \centering
  \schema[5.3cm]{scene}
  \fil{1}
  \astcue{
    \cue{une scène = un nom, une taille et une \cuekey{fonction} qui remplit le monde}
    \cue{\code{R} recommence la scène : on rappelle \code{build} sur un monde neuf}
    \cue{les forces sont déclarées \cuekey{avant} \code{world} : détruites \cuekey{après} lui}
    \cuewarn{les particules gardent des références vers ces forces jusqu'au bout}
    \cuenext{ce que fait \code{build} : créer des particules}
  }
\end{frame}

% ============================================================ 2. particules
\begin{frame}{Créer une particule}
  \framesubtitle{constructeurs délégués et liste d'initialisation}
  \centering
  \schema[5.3cm]{particules}
  \fil{2}
  \astcue{
    \cue{trois constructeurs, \cuekey{un seul fait le travail} : les autres délèguent}
    \cue{\code{radius} est \code{const} : impossible de l'affecter dans le corps}
    \cue{\code{static} : le compteur appartient à la classe, pas à un objet}
    \cuenext{où vivent ces particules ? dans le \code{World}}
  }
\end{frame}

% ============================================================ 3. world
\begin{frame}{Le \code{World} possède tout}
  \framesubtitle{des adresses qui ne bougent jamais}
  \centering
  \schema[5.3cm]{world}
  \fil{3}
  \astcue{
    \cue{un ressort garde \cuekey{deux pointeurs} vers ses particules}
    \cue{\code{vector} : \code{push\_back} peut tout déplacer $\to$ pointeurs pendants}
    \cue{\code{deque} : les éléments ne bougent jamais}
    \cue{règle des trois : copie \cuekey{interdite} plutôt que copie profonde, choix justifié}
    \cuewarn{c'était le bug de la première version : ressorts sans effet}
    \cuenext{le monde est prêt, on lance un pas : les forces}
  }
\end{frame}

% ============================================================ 4. forces
\begin{frame}{Une particule somme ses forces}
  \framesubtitle{un seul appel, quatre comportements}
  \centering
  \schema[5.3cm]{forces}
  \fil{4}
  \astcue{
    \cue{chaque particule tient une liste de \code{Force\&} : \cuekey{liaison dynamique}}
    \cue{\code{Spring} s'inscrit lui-même sur ses deux particules}
    \cue{\code{World} ne fait aucun \code{if} sur le type d'une force}
    \cuewarn{limite : un ressort est calculé deux fois, une par extrémité}
    \cuenext{zoom sur le ressort : la formule s'écrit avec \code{Vec2}}
  }
\end{frame}

% ============================================================ 4b. Vec2
\begin{frame}[fragile]{Le ressort s'écrit comme la formule}
  \framesubtitle{surcharge d'opérateurs dans \code{Vec2}}
  \begin{columns}[c,onlytextwidth]
    \begin{column}{0.66\textwidth}
      \centering
      {\large
      $\vec F_A = \big[\,\textcolor{astblue}{k\,(\lVert\vec d\rVert - L_0)}
        + \textcolor{timaorange}{\rho\,(\vec v_B-\vec v_A)\cdot\hat d}\,\big]\,
        \textcolor{phelmared}{\hat d}$}
      \par\vspace{3mm}
\begin{lstlisting}[escapeinside={(*@}{@*)},basicstyle=\ttfamily\tiny]
const Vec2 d = ends[1]->getPos() - ends[0]->getPos();
const float length = d.norm();
const Vec2 (*@\textcolor{phelmared}{direction = d / length}@*);
const float (*@\textcolor{timaorange}{speed = (ends[1]->getSpeed() - ends[0]->getSpeed())}@*)
                    (*@\textcolor{timaorange}{.dot(direction)}@*);
return direction * ((*@\textcolor{astblue}{k * (length - L0)}@*) + (*@\textcolor{timaorange}{rho * speed}@*));
\end{lstlisting}
      \hfill\filechip{physics/forces/springs.cpp, abrégé}
    \end{column}
    \begin{column}{0.31\textwidth}
      \begin{tikzpicture}[x=1cm,y=1cm]
        \node[cls=datagreen/datafill,text width=3.9cm,align=left] (v) at (0,2.6)
          {\clsname{Vec2}\\[0.4ex]
           {\scriptsize méthodes :\\\code{+ - * / += -= == !=}\\[0.5ex]
            fonction libre :\\\code{float * Vec2}\\[0.5ex]
            \code{dot}, \code{norm}, \code{normalized}}};
        \node[astcallout,text width=3.7cm,font=\scriptsize] at (0,0.35)
          {\code{const Vec2\&} en paramètre\\méthodes \code{const} : lecture seule};
      \end{tikzpicture}
    \end{column}
  \end{columns}
  \fil{4}
  \astcue{
    \cue{chaque couleur de la formule se retrouve \cuekey{telle quelle} dans le code}
    \cue{\code{k * x} marche grâce à la fonction libre \code{operator*(float, Vec2)}}
    \cue{passage par \code{const\&} : pas de copie, pas de modification}
    \cuenext{on a l'accélération : il faut avancer dans le temps}
  }
\end{frame}

% ============================================================ 5. integration
\begin{frame}{Changer d'intégrateur en direct}
  \framesubtitle{\code{enum class}, une fabrique et une interface}
  \centering
  \schema[5.3cm]{integration}
  \fil{5}
  \astcue{
    \cue{\code{enum class} : un type à part, pas un \code{int} déguisé}
    \cue{toute la décision est dans la \cuekey{fabrique}, hors de \code{World}}
    \cue{démo : appuyer sur \cuekey{I} pendant le tissu}
    \cuenext{lequel est le meilleur ? on a mesuré l'énergie}
  }
\end{frame}

% ============================================================ 5b. energie
% [fragile] + \shorthandoff : pgfplots lit les << ; >> que babel-french rend actifs
\begin{frame}[fragile]{Euler ou Verlet : l'énergie}
  \framesubtitle{écart à l'énergie initiale, 10\,s, sans amortissement}
  \centering\shorthandoff{;:}%
  \begin{tikzpicture}
    \begin{axis}[astplot,name=a1,title={Pendule},xlabel={$t$ (s)},ylabel={écart (\%)},
        xmin=0,xmax=10,ymin=-1.2,ymax=1.2,
        legend pos=south west]
      \addplot[timaorange] table[x=t,y expr={(\thisrow{euler}/58.86-1)*100},col sep=comma]
        {figures/data/energie_pendule.csv};
      \addplot[astblue,dash pattern=on 3pt off 1.5pt] table[x=t,y expr={(\thisrow{verlet}/58.86-1)*100},col sep=comma]
        {figures/data/energie_pendule.csv};
      \legend{Euler,Verlet}
    \end{axis}
    \begin{axis}[astplot,at={(a1.east)},anchor=west,xshift=1.4cm,title={Double pendule (chaotique)},
        xlabel={$t$ (s)},xmin=0,xmax=10,ymin=-4,ymax=4]
      \addplot[timaorange] table[x=t,y expr={(\thisrow{euler}/100.847-1)*100},col sep=comma]
        {figures/data/energie_double_pendule.csv};
      \addplot[astblue,dash pattern=on 3pt off 1.5pt] table[x=t,y expr={(\thisrow{verlet}/100.847-1)*100},col sep=comma]
        {figures/data/energie_double_pendule.csv};
    \end{axis}
  \end{tikzpicture}
  \par\vspace{1mm}
  \begin{tikzpicture}
    \node[cls=datagreen/datafill,font=\scriptsize,align=center]
      {critère du sujet tenu : \textbf{moins de 5\,\%} sur 10\,s avec Verlet (pendule : 0,7\,\%)\\[0.2ex]
       {\tiny\color{textgrey}Euler \emph{semi-implicite} conserve aussi bien l'énergie : les deux schémas sont symplectiques}};
  \end{tikzpicture}
  \fil{5}
  \astcue{
    \cue{mesuré par le programme \code{energie} : $\frac12 mv^2 + mgh + \frac12 k(\ell-L_0)^2$}
    \cue{les deux courbes se superposent : \cuekey{pas de dérive}}
    \cuewarn{on attendait Euler moins bon : c'est vrai pour l'Euler \emph{explicite}, pas le semi-implicite}
    \cuenext{après l'intégration : on remet les particules dans la boîte}
  }
\end{frame}

% ============================================================ 6. contraintes
\begin{frame}{Les contraintes}
  \framesubtitle{bords du monde, rebond et points épinglés}
  \begin{columns}[c,onlytextwidth]
    \begin{column}{0.70\textwidth}
      \centering\schema[5.3cm]{contraintes}
    \end{column}
    \begin{column}{0.28\textwidth}
      \centering\scenecard{figures/scenes/scene_10.png}{Pluie de balles}{2.2cm}
    \end{column}
  \end{columns}
  \fil{6}
  \astcue{
    \cue{le monde est une \cuekey{boîte} : 4 bords au lieu d'un simple sol}
    \cue{coefficient de restitution $e$ : chaque balle a le sien}
    \cue{Verlet déduit la vitesse des positions : il faut corriger son historique}
    \cuenext{le pas est fini, on dessine}
  }
\end{frame}

% ============================================================ 7. dessin / diamant
\begin{frame}{Dessiner : le diamant}
  \framesubtitle{\code{Spring} est une \code{Force} et un \code{Drawable}}
  \centering
  \schema[5.3cm]{diamant}
  \fil{7}
  \astcue{
    \cue{\code{Force} et \code{Drawable} héritent \cuekey{virtuellement} d'\code{Entity}}
    \cue{mesure \code{sizeof} : 112 contre 144 octets, et un seul id}
    \cue{le moteur appelle \code{Renderer} : SFML reste une \cuekey{boîte noire}}
    \cuenext{dernière étape : l'utilisateur attrape une particule}
  }
\end{frame}

% ============================================================ 8. souris
\begin{frame}[fragile]{Attraper une particule à la souris}
  \framesubtitle{une lambda et un \code{dynamic\_cast}}
  \begin{columns}[c,onlytextwidth]
    \begin{column}{0.56\textwidth}
\begin{lstlisting}
Particle *World::pickParticle(const Vec2 &at,
                              float max_dist) {
  Particle *best = nullptr;
  float best_dist = max_dist;
  // la lambda capture best et best_dist
  auto closer = [&](Entity *e) {
    if (Particle *p = dynamic_cast<Particle *>(e)) {
      const float d = (p->getPos() - at).norm();
      if (d < best_dist) { best_dist = d; best = p; }
    }
  };
  std::vector<Entity *> all = entities();
  std::for_each(all.begin(), all.end(), closer);
  return best;
}
\end{lstlisting}
      \hfill\filechip{physics/world.cpp}
    \end{column}
    \begin{column}{0.40\textwidth}
      \centering
      \schema{souris}
    \end{column}
  \end{columns}
  \fil{8}
  \astcue{
    \cue{\code{dynamic\_cast} renvoie \code{nullptr} pour un ressort}
    \cue{la lambda capture par référence : \cuekey{[\&]}}
    \cue{relâcher = lancer : la vitesse de la souris est rendue à la particule}
    \cuenext{ce que tout ça donne : les scènes}
  }
\end{frame}
       % un pas de bout en bout : scene, particules,
                                   % World, forces, Vec2, integration, energie,
                                   % contraintes, dessin (diamant), souris
% !TEX root = ../main.tex
% Partie 3 : ce que le moteur produit (extension, scenes), recapitulatif des
% notions et bilan.

\section{Résultats}

% ============================================================ extension : vent
\begin{frame}[fragile]{L'extension : un vent qui varie}
  \framesubtitle{une force dont l'intensité est une lambda du temps}
  \begin{columns}[c,onlytextwidth]
    \begin{column}{0.45\textwidth}
      \centering
      \scenecard{figures/scenes/scene_07.png}{Rideau au vent}{3.9cm}
    \end{column}
    \begin{column}{0.53\textwidth}
      \centering\shorthandoff{;:}%
      \begin{tikzpicture}
        \begin{axis}[astplot,width=7.4cm,height=3.2cm,xlabel={$t$ (s)},ylabel={$w_x$ (m/s)},
            xmin=0,xmax=20,ymin=-2,ymax=14,ytick={0,6,12},samples=300,domain=0:20]
          \addplot[violet!70!black,fill=sinkfill,fill opacity=0.5] {6+5*sin(deg(0.7*x))+2*sin(deg(2.3*x+1))} \closedcycle;
          \addplot[sinkpurple,line width=1.2pt] {6+5*sin(deg(0.7*x))+2*sin(deg(2.3*x+1))};
        \end{axis}
      \end{tikzpicture}
\begin{lstlisting}[basicstyle=\ttfamily\tiny]
Wind &wind = s.addForce<Wind>(
    [](float t) {   // rafales : deux sinus
      float gust = 6 + 5*sin(0.7f*t) + 2*sin(2.3f*t + 1);
      return Vec2(gust, 0.8f * sin(1.3f*t));
    },
    s.time, 0.035f);           // F = c (w(t) - v)
\end{lstlisting}
    \end{column}
  \end{columns}
  \vspace{1mm}
  \centering
  \begin{tikzpicture}
    \node[cls=datagreen/datafill,font=\scriptsize]
      {une classe \code{Wind} + une lambda \astsep{} \code{World}, \code{Integrator} et les autres forces \textbf{inchangés}};
  \end{tikzpicture}
  \astcue{
    \cue{extension du jalon 4 : \cuekey{vent variable}}
    \cue{\code{std::function<Vec2(float)>} : n'importe quel profil de vent}
    \cue{la force lit l'horloge de la scène, elle ne modifie jamais \code{World}}
    \cuenext{toutes les scènes}
  }
\end{frame}

% ============================================================ galerie
\begin{frame}{Treize scènes, un seul moteur}
  \framesubtitle{une scène = une fonction qui remplit un \code{World}}
  \centering
  \setlength{\tabcolsep}{1.5mm}
  \begin{tabular}{ccc}
    \scenecard{figures/scenes/scene_04.png}{Pont de corde}{2.6cm} &
    \scenecard{figures/scenes/scene_12.png}{Toile d'araignée}{2.6cm} &
    \scenecard{figures/scenes/scene_09.png}{Oscillateurs}{2.6cm} \\[3mm]
    \scenecard{figures/scenes/scene_06.png}{Blob}{2.6cm} &
    \scenecard{figures/scenes/scene_08.png}{Gelée}{2.6cm} &
    \scenecard{figures/scenes/scene_11.png}{Apesanteur}{2.6cm} \\
  \end{tabular}
  \astcue{
    \cue{ajouter une scène = écrire une fonction, \cuekey{sans toucher au moteur}}
    \cue{oscillateurs : 6 raideurs, de 5 à 160 N/m, même masse}
    \cue{apesanteur : $g = 0$, un anneau qui tourne sur lui-même}
    \cuenext{démo en direct si le temps le permet}
  }
\end{frame}

% ============================================================ recapitulatif
\begin{frame}{Les notions du cours, où les trouver}
  \framesubtitle{du plus simple au plus avancé}
  \centering
  \begin{tikzpicture}[x=1cm,y=1cm]
    \foreach \i/\c/\f/\t/\b/\file in {
      0/datagreen/datafill/Vec2/{opérateurs\\passage par \code{const\&}\\méthodes \code{const}}/vec2.hpp,
      1/timagreen/timagreen!12/Constructeurs/{délégués\\liste d'initialisation\\\code{const}, \code{static}}/particles.cpp,
      2/memborder/memfill/Polymorphisme/{\code{virtual}, \code{override}\\\code{enum class}\\fabrique}/force.hpp,
      3/sinkpurple/sinkfill/Mémoire/{\code{deque}\\\code{unique\_ptr}\\copie \code{= delete}}/world.hpp,
      4/ctrlred/ctrlfill/{Héritage\\multiple}/{diamant virtuel\\\code{dynamic\_cast}\\lambdas, template}/springs.hpp}
    {
      \node[cls=\c/\f,anchor=south west,minimum width=2.75cm,minimum height=2.3cm,
        text width=2.5cm,align=left,inner sep=1.3mm] (st\i) at (\i*2.95,\i*0.85)
        {{\bfseries\footnotesize\t}\\[0.8ex]{\scriptsize\b}};
      \node[astnum] at (st\i.north west) {\the\numexpr\i+1\relax};
      \node[anchor=south east,inner sep=0.6mm] at (st\i.north east) {\filechip{\file}};
    }
    \draw[-{Latex[length=2.4mm]},line width=1.4pt,phelmared,rounded corners=1mm]
      (0,-0.15) -- (2.85,-0.15) -- (2.85,0.7) -- (5.8,0.7) -- (5.8,1.55) -- (8.75,1.55)
      -- (8.75,2.4) -- (11.7,2.4) -- (11.7,3.25) -- (14.65,3.25) -- (14.65,4.6);
    \node[anchor=north east,font=\tiny\bfseries,text=phelmared,align=right] at (14.55,0.5)
      {plus abstrait,\\plus puissant};
  \end{tikzpicture}
  \astcue{
    \cue{le tableau du sujet en une image : \cuekey{chaque notion a son fichier}}
    \cuenext{bilan}
  }
\end{frame}

% ============================================================ bilan
\begin{frame}{Bilan}
  \framesubtitle{ce que le moteur fait, et ses limites}
  \centering
  \begin{tikzpicture}[x=1cm,y=1cm]
    \node[cls=datagreen/datafill,text width=6.3cm,align=left,minimum height=4.6cm,anchor=north] (ok) at (3.5,5.0)
      {{\bfseries\small\color{datagreen!50!black}\faCheck\enspace Ce qui marche}\\[1.2ex]
       {\scriptsize
        \textbf{13 scènes}, du pendule au tissu de 400 particules\\[0.7ex]
        Euler ou Verlet, \textbf{changés en direct}\\[0.7ex]
        \code{physics} compilée en \code{.a}, \textbf{sans SFML}\\[0.7ex]
        une force de plus = \textbf{une classe} (le vent)\\[0.7ex]
        énergie : moins de 1\,\% d'écart sur le pendule\\[0.7ex]
        \code{-Wall -Wextra -Werror} : \textbf{zéro avertissement}\\[0.7ex]
        test sans affichage : 13 scènes $\times$ 2 intégrateurs}};
    \node[cls=ctrlred/ctrlfill,text width=6.3cm,align=left,minimum height=4.6cm,anchor=north] (ko) at (11.1,5.0)
      {{\bfseries\small\color{ctrlred}\faExclamationTriangle\enspace Limites connues}\\[1.2ex]
       {\scriptsize
        pas de collision \textbf{entre} particules, seulement avec les bords\\[0.7ex]
        un ressort est calculé \textbf{deux fois}, une par extrémité\\[0.7ex]
        Verlet : vitesse estimée par différence arrière\\[0.7ex]
        forces attachées aux particules, pas au \code{World}\\[0.7ex]
        pas de rendu SVG ni de \code{namespace}}};
  \end{tikzpicture}
  \astcue{
    \cue{insister sur l'\cuekey{extensibilité} : vent ajouté sans toucher au moteur}
    \cue{limites : les dire nous-mêmes avant qu'on nous les demande}
    \cuenext{questions, puis démo}
  }
\end{frame}

% ============================================================ fin
\begin{frame}[plain]
  \centering
  \vspace*{\fill}
  \ressortsemblem[1.3]\par\vspace{6mm}
  {\LARGE\bfseries\color{inpnavy}Merci}\par\vspace{2mm}
  {\color{phelmared}\rule{2.6cm}{1.2pt}}\par\vspace{3mm}
  {\normalsize\color{textgrey}Questions \astsep{} démo}\par\vspace{5mm}
  {\scriptsize\color{textgrey}\code{github.com/Audr1c/POOProject}\astsep branche \code{presentation}}
  \vspace*{\fill}
  \astcue{
    \cue{lancer \code{./poop} : tissu, \cuekey{I} pour changer d'intégrateur, rideau au vent}
    \cue{attraper, lancer, épingler (clic droit) une particule}
  }
\end{frame}
      % vent, scenes, notions, bilan
\end{document}
